\section{衝突演算子と数値安定性}
\label{sec:cumulant}

単純なBGK（単一緩和時間）衝突演算子は，高レイノルズ数・低粘性域（\(\tau\to1/2\)）で数値不安定になりやすい。本番計算が要求する \(Re\sim2\times10^5\) では \(\tau-1/2\) が\(10^{-5}\)のオーダーとなるため，BGKでは実用にならない。本プロジェクトでは，安定性を確保するために\textbf{中心モーメント（キュムラント緩和）衝突演算子}を実装した（\texttt{src/core/central\_moments.jl}）。

\subsection{中心モーメント衝突演算子}

中心モーメント衝突演算子は，流体とともに動く座標系でモーメントを取り，その意味ごとに異なる緩和率を与える。

\begin{itemize}
	\item 2次モーメントの偏差成分（せん断応力を担う）は，動粘性を決める緩和率 \(\omega=1/\tau\) で緩和する。
	\item その対角和（トレース，音響＝体積粘性モード）は独立の緩和率 \(\omega_b\approx1\) で緩和する。これを2次の偏差成分と分離することが，安定性向上の主因である。
	\item 3次以上のモーメントは，キュムラントを0へ緩和する。これは「平衡値へ固定する」のとは異なり，\emph{その時点の2次モーメントが含意する値}（Isserlisの定理によるガウス閉包）へ緩和することを意味する。
	\item \label{item:wall-position}奇数次と偶数次で緩和率を分ける（\(\omega_{odd}\) / \(\omega_{higher}\)，既定はいずれも \(1.0\)）。奇数次の緩和率はバウンスバックが実効壁位置として読む\(\omega^-\)であり，TRT（Two-Relaxation-Time）の magic parameter \(\Lambda\) を自由に選べるようにする唯一の自由度である。実測では既定の\(1.0\)が最適に近く，\(\Lambda=3/16\)（TRTが厳密に壁位置を格子中間に置くとされる値）に合わせるとむしろ壁位置は悪化する（\secref{validation}）。偶数次の緩和率を奇数次と一緒に下げると，\(Re=100\)で200ステップ以内に発散するため，緩和率を分けること自体が安定性の前提になっている。
\end{itemize}

\paragraph{TRTのmagic parameterと実効壁位置}
\label{sec:wall-position}
補間バウンスバックは，格子の幾何がある位置に壁を置いていても，数値的に「見える」壁の位置（実効壁位置）はそこから系統的にずれうる。Two-Relaxation-Time（TRT）の理論では，このずれはmagic parameter
\begin{align}
	\Lambda = \left(\frac{1}{\omega^+}-\frac{1}{2}\right)\left(\frac{1}{\omega^-}-\frac{1}{2}\right)
	\label{eq:magic-parameter}
\end{align}
で特徴づけられ，\(\Lambda=3/16\)のとき壁は格子の中間に置かれるとされる。\(\omega^+\)はせん断粘性を担う偶数次の緩和率（\(1/\omega^+-1/2=\tau-1/2\)であり自由に選べない）で，\(\omega^-\)が上記の奇数次緩和率\(\omega_{odd}\)にあたる。BGKでは両者が同一であるため\(\Lambda=(\tau-1/2)^2\)に固定されてしまうが，中心モーメント演算子で緩和率を分離したことにより\(\Lambda\)を独立に設定できるようになる。この自由度が壁位置に関する誤差の原因かどうかは\secref{validation-highre}で直接検証する。

D3Q27の離散速度セットが \(\{-1,0,1\}^3\) のテンソル積であることを利用し，モーメントへの変換は3回の1次元変換（各3点を0/1/2次モーメントへ写す）に分解され，閉形式で逆変換できる。これにより演算コストを抑えつつ，27方向すべての情報を扱える。

\subsection{安定性と精度の実測比較}

BGK（D3Q19）と中心モーメント演算子（D3Q27）を，二重周期せん断層とTaylor-Green渦で比較した（\texttt{test/test\_central\_moments.jl}）。

\begin{table}[H]
	\centering
	\caption{BGKと中心モーメント演算子の安定性・精度比較。}
	\label{tab:collision-comparison}
	\begin{tabular}{lcc}
		\Hline{1.5pt}
		項目 & BGK (D3Q19) & 中心モーメント (D3Q27) \\
		\hline
		二重周期せん断層が安定な下限 \(\tau\) & 0.502（\(Re\approx9{,}600\)） & \textbf{0.5001}（\(Re\approx192{,}000\)） \\
		Taylor-Green相対L2誤差（\(n=16/32/64\)） & 5.40e-3 / 1.33e-3 / 3.45e-4 & \textbf{1.63e-3} / \textbf{4.10e-4} / \textbf{1.25e-4} \\
		空間収束次数 & 2.03 / 1.94 & 1.99 / 1.72 \\
		\hline
	\end{tabular}
\end{table}

中心モーメント演算子はBGKが\(\tau=0.501\)で発散する条件で\(\tau=0.5001\)まで安定であり，全格子でBGKより高精度である。安定性向上の主因は体積粘性モードの分離であり，これを行わずに全2次モーメントを\(\omega\)で一括緩和する実装では，せん断層の安定限界はBGKとほぼ変わらないことを確認している。

\paragraph{既知の限界}
離散速度が \(c\in\{-1,0,1\}\) を取るため \(c^3=c\) となり，3次中心モーメントが \(-\rho u^3\) に固定される（Maxwell分布であれば0）。これは速度セット由来の誤差であり，どの衝突演算子を用いても除去できない。一様流を重ねた場合のガリレイ不変性の誤差は，中心モーメント演算子でもBGKと同程度残る。Geierらのキュムラント法\cite{geier2015cumulant}は，これを速度勾配から構成した補正項で打ち消すが，本実装では未対応である。
